\documentclass[a4paper,12pt]{report}
\usepackage{algorithmic}
\usepackage[linesnumbered,ruled,vlined]{algorithm2e}
\usepackage[margin=2cm]{geometry}
\usepackage[utf8]{inputenc}
\usepackage{listings} 
\usepackage{graphicx} 
\usepackage{color}
\usepackage{xcolor}
\usepackage{hyperref}
\usepackage{verbatim}
%\usepackage{mdframed}

\newcommand{\currentdata}{14 February 2015}
\newtheorem{example}{Example}

\begin{document}
\vspace{-5cm}
\begin{center}
Department of Computer Science\\
Technical University of Cluj-Napoca\\
\includegraphics[width=10cm]{fig/footer}
\end{center}
\vspace{1cm}
%\maketitle
\begin{center}
\begin{Large}
 \textbf{Artificial Intelligence}\\
\end{Large}
\textit{Laboratory activity}\\
\vspace{3cm}
Meze Razvan-Gabriel \& Nedelcu Ioan-Andrei\\
Group: 2\\
Email: tatarazvi@gmail.com anedelcu0226@gmail.com\\
\vspace{12cm}
Teaching Assistant: Adrian Groza\\
Adrian.Groza@cs.utcluj.ro\\
\vspace{1cm}
\includegraphics[width=10cm]{fig/footer}
\end{center}

\tableofcontents

% %\chapter{Rules and policies}
\newpage
\paragraph{Lab organisation.} 

\begin{enumerate}
 \item Laboratory work is 25\% from the final grade. 
 \item There are three deliverables in total: 1. Search, 2. Logic, 3. Planning.
%\item Tables~\ref{tab:first} and~\ref{tab:second} also bears out nine deadlines that you have to meet for achieving the points. They correspond to the following weeks: $W_2, W_4, W_5, W_6, W_{10}, W_{11}, W_{12}, W_{13}, W_{14}$. After each deadline you will get (or not) the achievable points. 
\item Before each deadline, you have to send your work (latex documentation/code) at moodle.cs.utcluj.ro

\item We use Linux and Latex
\item Plagiarism: Don't be a cheater! Cheating affects your colleagues, scholarships and a lot more.
\end{enumerate}
\begin{table}
\begin{center}
\caption{Lab scheduling}
\begin{tabular}{p{11cm}p{1.7cm}}\hline
\textbf{Activity} & \textbf{Deadline}  \\ \hline
{\it Searching agents, Linux, Latex, Python, Pacman}  & $W_1$ \\ %\hline
{\it Uninformed search} & $W_2$  \\ %\hline
{\it Informed Search} & $W_3$  \\ %\hline
{\it Adversarial search} & $W_4$ \\ %\hline
{\it Propositional logic} & $W_5$ \\ %\hline
{\it First order logic} & $W_6$ \\ %\hline
{\it Inference in first order logic} & $W_7$ \\ %\hline
{\it Knowledge representation in first order logic} & $W_8$ \\ %\hline
{\it Classical planning} & $W_9$ \\ %\hline
{\it Contingent, conformant and probabilistic planning} & $W_{10}$ \\ % \hline
{\it Multi-agent planing} & $W_{11}$ \\ %\hline
{\it Modelling planning domains} & $W_{12}$ \\ %\hline
{\it Planning with event calculus} & $W_{14}$ \\ \hline 
\end{tabular}
\end{center}
\label{tab:first}
\end{table}


\begin{comment}
\begin{center}
\begin{tabular}{ll}
Class: & Introducere in Inteligenta Artificiala\\
Enrollment key: & Iia2017-2018\\
\end{tabular}

\end{center}


\item {\it Laptop policy}: you can use your own laptop as long you have Linux. One goal of the laboratory is to increase your competency in Linux. It is \textbf{your} task to set static IPs:
\begin{center}
\begin{tabular}{ll}
IP: & 192.168.1.51\footnote{Or 52-59: avoid introducing the same IP for two different stations.}\\ 
MASK:&  255.255.255.0\\
GATEWAY:&  192.168.1.2\\
DNS:& 192.168.1.2\\
PROXY & 192.168.1.2:3128\\
\end{tabular}

Wifi:
\begin{tabular}{ll}
Network: &isg\\
Password: &inteligentaartificiala\\
\end{tabular}
\end{center}

\item {\it Group change policy.} Maximum number of students in a class is 14.

\item {\it For students repeating the class}: A discussion for validating the previous grade is mandatory in the first week. I usually have no problem to validate your previous grades, as long you request this in the first week. 
Failing to do so, leads to the grade 1 for the laboratory work in the current semester.
\end{comment}



\begin{comment}
\paragraph{Grading.}
Assessment aims to measure your knowledge and skills needed to function in realistic AI-related tasks.
Assessment is based on your written report explaining the nature of the project, findings, and recommendations. 
Meeting the deadlines is also important.
Your report is comparable to ones you would write if you were a consultant reporting to a client.

Grade inflation makes difficult to distinguish between students.
It also discourages the best students to do their best.
In my quest for ``optimal ranking of the students``, I do not use the following heuristics:

\begin{itemize}
 \item "He worked hard at the project``. 
Our society do not like anymore individuals that are $trying$, but individual that $do$ stuff.
Such heuristic is not admissible in education, except the primary school.

\item "I knew he could do much better". 
Such a heuristic is not admissible because it does not encourage you to spread yourself.
\end{itemize}

\begin{itemize}

\item 7 means that you: 
i) constantly worked during classes, 
 ii) you proved competent to use the tool and its expressivity for a realistic scenario, 
iii) you understood theoretical concepts on which the tool rely on.

\item 8, 9 mean that your code is large enough and the results proved by your experiments are significant.

 \item 10 means that you did very impressive work or more efficient 
that I expected or handled a lot of special cases for realistic scenarios.


\item 5 means that you managed to develop something of your own, functional, 
with your own piece of code substantially different from the examples available.

\item You obtain less than 5 in one of the following situations: 

\begin{enumerate}
\item few code written by yourself.
\item too much similarity with the provided examples.
\item non-seriosity (i.e. re-current late at classes, playing games, worked for other disciplines, 
poor/unprofessional documentation of your work, etc.)\footnote{Consider non-seriosity as a immutable boolean value that 
is unconsciously activated in my brain when one of the above conditions occurs for the first time.}.
\end{enumerate}

\item You get 2 if you present the project but fail to submit the documentation or code.
You get 1 if you do not present your project before the deadline.
You get 0 for any line of code taken from other parts that appear in section {\it My own code.} 
For information on TUCN's regulations on plagiarism do consult the active norms.

If your grade is 0, 1, or 2, you do not satisfy 
the preconditions for participating to the written exam. 
The only possibility to increase your laboratory grade is to take another project in the next year, 
at the same class, and to make all the steps again.   
\end{itemize}

However, don't forget that focus is on learning, not on grading.

{\it Using Latex in your documentation.} You have to show some competency on writing documentation in Latex. 
For instance, you have to employ various latex elements: lists, citations, footnotes, verbatim, maths, code, etc. 



\paragraph{Plagiarism.} Most of you consider plagiarism only a minor form of cheating.
This is far from accurate. Plagiarism is passing off the work of others as your own to gain unfair advantage. 

During your project presentation and documentation, 
I must not be left with doubts on which parts of your project are your work or not.  
Always identify both: 
1) who you worked with and 
2) where you got your part of the code or solution.
You should sign the declaration of originality. 

Describe clearly the starting point of your solution.
List explicitly any code re-used in your project. 
List explicitly any help (including debugging help, design discussions) provided by others (including colleagues or teaching assistant). 
Keep in mind that it is your own project and not the teaching assistant's project. 
Learning by collaborating does remain an effective method. 
You can use it, but don't forget to mention any kind of support. 
Learning by exploiting various knowledge-bases developed by your elder colleagues remain also an effective method for ''learning by example''. 
When comparing samples of good and poor assignments submitted by your colleagues in earlier years try to identify which is better and why. 
You can use this repository of previous assignments, but don't forget to mention any kind of inspiration source. 

The assignment is designed to be individual and to pose you some difficulties 
(both technological and scientific) for which you should identify a working solution by the end of the semester.  
Each semester, a distinct AI tool is assigned to two students. 
Your are encouraged to collaborate, especially during the the installation and example understanding phases ($W_1$-$W_4$). The quicker you get throughout these preparatory stages, the more time you have for your own project.


\paragraph{Class attendance.} 
I expect active participation at all activities. %We are all grown-ups, when and whether you attend lecture is up to you\footnote{However, you should be aware that when signing the study contract, you have some moral obligations towards the people that are paying your classes through the Ministry of Education budget.}. 
Keep in mind the exam can include any topic that was covered during class, explained on the board, 
or which emerged from discussions among participants. 
Missing lab assignments or midterm leads to minimum grade for that part.
You are free to manage your laboratory classes - meaning that you can submit the project earlier -  as long as you meet all the constraints and deadlines.
%However, it is mandatory to participate at the final public presentation of your project. 
\end{comment}

%\chapter{Laboratory works}

\chapter{Logics}
\section{Description of the Input File}

The prisoner's dilemma game is described using logical formulas. The game consists of two prisoners, and each prisoner has the choice to confess (C1 and C2) or not confess (-C1 and -C2). Each choice leads to a specific outcome for each prisoner, represented by the variables LS1, MS1, SS1, NS1 for prisoner 1 and LS2, MS2, SS2, NS2 for prisoner 2, which stand for large sentence, medium sentence, small sentence, and no sentence, respectively.

The file also includes formulas that describe the alternative strategies that each prisoner can choose, represented by the variables alternative\_strategy\_C1, alternative\_strategy\_not\_C1, alternative\_strategy\_C2, and alternative\_strategy\_not\_C2, which stand for prisoner 1 switching to confessing or not confessing, and prisoner 2 switching to confessing or not confessing, respectively.

For each alternative strategy, there is an associated alternative sentence for each prisoner, represented by the variables LAS1, MAS1, SAS1, NAS1 for prisoner 1 and LAS2, MAS2, SAS2, NAS2 for prisoner 2.

The file also includes formulas that describe the benefits that each prisoner can gain by switching strategies, represented by the variables P1\_change\_benefit and P2\_change\_benefit, which stand for prisoner 1 and prisoner 2 benefiting, respectively, if they switch strategies from the current state and the other prisoner keeps the same strategy.

Finally, the file includes a formula for the concept of Nash equilibrium, represented by the variable nash\_equilibrium, which indicates that the current state is a Nash equilibrium if no prisoner gets a smaller sentence if he switches strategies and the other prisoner keeps the same strategy.

The goal of the MACE4 input file is to find all the models in which a Nash equilibrium has been reached, in this case, a unique model in which both prisoners confess.

The input file for MACE4 describes the prisoner's dilemma game, in which two prisoners must decide whether to confess or not confess. The game is described using logical formulas.

\subsection{Variables}

The following variables are used in the input file:

\begin{itemize}
\item $C1$: prisoner 1 confesses
\item $C2$: prisoner 2 confesses
\item $LS1$: prisoner 1 gets a large sentence
\item $MS1$: prisoner 1 gets a medium sentence
\item $SS1$: prisoner 1 gets a small sentence
\item $NS1$: prisoner 1 gets no sentence
\item $LS2$: prisoner 2 gets a large sentence
\item $MS2$: prisoner 2 gets a medium sentence
\item $SS2$: prisoner 2 gets a small sentence
\item $NS2$: prisoner 2 gets no sentence
\end{itemize}

\subsection{Logical Formulas}

The input file includes the following logical formulas:

\begin{itemize}
\item If both prisoners do not confess, then both prisoners get small sentences: \
$(-C1 \land -C2) \leftrightarrow (-LS1 \land -LS2 \land -MS1 \land -MS2 \land SS1 \land SS2 \land -NS1 \land -NS2)$
\item If prisoner 1 does not confess and prisoner 2 confesses, then prisoner 1 gets a large sentence and prisoner 2 gets no sentence: \
$(-C1 \land C2) \leftrightarrow (LS1 \land -LS2 \land -MS1 \land -MS2 \land -SS1 \land -SS2 \land -NS1 \land NS2)$
\item If prisoner 1 confesses and prisoner 2 does not confess, then prisoner 1 gets no sentence and prisoner 2 gets a large sentence: \
$(C1 \land -C2) \leftrightarrow (-LS1 \land LS2 \land -MS1 \land -MS2 \land -SS1 \land -SS2 \land NS1 \land -NS2)$
\item If both prisoners confess, then both prisoners get medium sentences: \
$(C1 \land C2) \leftrightarrow (-LS1 \land -LS2 \land MS1 \land MS2 \land -SS1 \land -SS2 \land -NS1 \land -NS2)$
\item $alternative\_strategy\_C1$: prisoner 1 can switch his decision to confessing, from the current state
\item $alternative\_strategy\_not\_C1$: prisoner 1 can switch his decision to not confessing, from the current state
\item $alternative\_strategy\_C2$: prisoner 2 can switch his decision to confessing, from the current state
\item $alternative\_strategy\_not\_C2$: prisoner 2 can switch his decision to not confessing, from the current state
\item $LAS1$: prisoner 1 gets a large sentence if he switches to an alternative strategy
\item $MAS1$: prisoner 1 gets a medium sentence if he switches to an alternative strategy
\item $SAS1$: prisoner 1 gets a small sentence if he switches to an alternative strategy
\item $NAS1$: prisoner 1 gets no sentence if he switches to an alternative strategy
\item $LAS2$: prisoner 2 gets a large sentence if he switches to an alternative strategy
\item $MAS2$: prisoner 2 gets a medium sentence if he switches to an alternative strategy
\item $SAS2$: prisoner 2 gets a small sentence if he switches to an alternative strategy
\item $NAS2$: prisoner 2 gets no sentence if he switches to an alternative strategy
\end{itemize}



The input file includes the following logical formulas, starting from the line "% each pair of desicions...":

\begin{itemize}
\item Each strategy has an alternative:
\begin{itemize}
\item $C1 \leftrightarrow alternative\_strategy\_not\_C1$
\item $C1 \leftrightarrow -alternative\_strategy\_C1$
\item $-C1 \leftrightarrow alternative\_strategy\_C1$
\item $-C1 \leftrightarrow -alternative\_strategy\_not\_C1$
\item $C2 \leftrightarrow alternative\_strategy\_not\_C2$
\item $C2 \leftrightarrow -alternative\_strategy\_C2$
\item $-C2 \leftrightarrow alternative\_strategy\_C2$
\item $-C2 \leftrightarrow -alternative\_strategy\_not\_C2$
\end{itemize}
\item Each alternative strategy, if applied, for each prisoner, has an alternative sentence associated with it (considering that the other prisoner has the same strategy):
\begin{itemize}
\item $alternative\_strategy\_C1 \land C2 \leftrightarrow -LAS1 \land MAS1 \land -SAS1 \land -NAS1$
\item $alternative\_strategy\_not$
\item $alternative\_strategy\_not\_C1 \land C2 \leftrightarrow LAS1 \land -MAS1 \land -SAS1 \land -NAS1$
\item $alternative\_strategy\_C1 \land -C2 \leftrightarrow -LAS1 \land -MAS1 \land -SAS1 \land NAS1$
\item $alternative\_strategy\_not\_C1 \land -C2 \leftrightarrow -LAS1 \land -MAS1 \land SAS1 \land -NAS1$
\item $C1 \land alternative\_strategy\_C2 \leftrightarrow -LAS2 \land MAS2 \land -SAS2 \land -NAS2$
\item $C1 \land alternative\_strategy\_not\_C2 \leftrightarrow LAS2 \land -MAS2 \land -SAS2 \land -NAS2$
\item $-C1 \land alternative\_strategy\_C2 \leftrightarrow -LAS2 \land -MAS2 \land -SAS2 \land NAS2$
\item $-C1 \land alternative\_strategy\_not\_C2 \leftrightarrow -LAS2 \land -MAS2 \land SAS2 \land -NAS2$
\end{itemize}
\end{itemize}

\begin{itemize}
\item $P1\_change\_benefit$: prisoner 1 benefits if he switches strategies from current state and the other prisoner keeps the same strategy (prisoner 1 gets a smaller sentence if he switches)
\item $P2\_change\_benefit$: prisoner 2 benefits if he switches strategies from current state and the other prisoner keeps the same strategy (prisoner 2 gets a smaller sentence if he switches)
\end{itemize}

\begin{itemize}
\item For each alternative strategy, there is a condition for which it is advantageous for the prisoner to change strategies:
\begin{itemize}
\item $alternative\_strategy\_C1 \land C2 \land -P1\_change\_benefit$
\item $alternative\_strategy\_not\_C1 \land C2 \land P1\_change\_benefit$
\item $alternative\_strategy\_C1 \land -C2 \land -P1\_change\_benefit$
\item $alternative\_strategy\_not_C1 \land -C2 \land P1\_change\_benefit$
\item $C1 \land alternative\_strategy\_C2 \land -P2\_change\_benefit$
\item $C1 \land alternative\_strategy\_not\_C2 \land P2\_change\_benefit$
\item $-C1 \land alternative\_strategy\_C2 \land -P2\_change\_benefit$
\item $-C1 \land alternative\_strategy\_not_C2 \land P2\_change\_benefit$
\end{itemize}
\end{itemize}

\begin{itemize}
\item For prisoner 1, it is advantageous to change strategies if the alternative sentence is no sentence or a small sentence:
\begin{align*}
& (NAS1 \land LS1) \lor (NAS1 \land MS1) \lor (NAS1 \land SS1) \lor (SAS1 \land LS1) \
& \quad \lor (SAS1 \land MS1) \lor (MAS1 \land LS1) \
& \leftrightarrow P1\_change\_benefit
\end{align*}
\item For prisoner 1, it is not advantageous to change strategies if the alternative sentence is a large sentence or a medium sentence:
\begin{align*}
& (LAS1 \land NS1) \lor (MAS1 \land NS1) \lor (SAS1 \land NS1) \lor (LAS1 \land SS1) \
& \quad \lor (MAS1 \land SS1) \lor (LAS1 \land MS1) \
& \leftrightarrow -P1\_change\_benefit
\end{align*}
\item For prisoner 2, it is advantageous to change strategies if the alternative sentence is no sentence or a small sentence:
\begin{align*}
& (NAS2 \land LS2) \lor (NAS2 \land MS2) \lor (NAS2 \land SS2) \lor (SAS2 \land LS2) \
& \quad \lor (SAS2 \land MS2) \lor (MAS2 \land LS2) \
& \leftrightarrow P2\_change\_benefit
\end{align*}
\item For prisoner 2, it is not advantageous to change strategies if the alternative sentence is a large sentence or a medium sentence:

\begin{align*}
& (LAS2 \land NS2) \lor (MAS2 \land NS2) \lor (SAS2 \land NS2) \lor (LAS2 \land SS2) \
& \quad
\end{align*}
\item For prisoner 2, it is not advantageous to change strategies if the alternative sentence is a large sentence or a medium sentence:
\begin{align*}
& (LAS2 \land NS2) \lor (MAS2 \land NS2) \lor (SAS2 \land NS2) \lor (LAS2 \land SS2) \
& \quad \lor (MAS2 \land SS2) \lor (LAS2 \land MS2) \
& \leftrightarrow -P2\_change\_benefit
\end{align*}

\item A Nash equilibrium has been reached if both prisoners do not benefit from changing strategies: \
$nash\_equilibrium \leftrightarrow (-P1\_change\_benefit \land -P2\_change\_benefit)$
\end{itemize}

\subsection{Goals}

The input file includes the following goal:

\begin{itemize}
\item Find all states that represent a Nash equilibrium: \
$-nash\_equilibrium$
\end{itemize}

\bibliographystyle{plain}
\bibliography{is}

\appendix

\chapter{Your original code}
Don't be a cheater! Cheating affects your colleagues, scholarships and a lot more.
This section should contain only code developed by you, without any line re-used from other sources. 
This section helps me to correctly evaluate your amount of work and results obtained. 


\vspace{2cm}
\begin{center}
Intelligent Systems Group\\
\includegraphics[width=10cm]{fig/footer}
\end{center}



\end{document}
